% ECONOMICS_PROBLEM: {"id": "L41-333", "title": "Dynamic competition, innovation and mergers", "jel_code": "L41", "source_id": "OP-0072"}
% FIRST_POSTED: 2026-10-09
% ATTEMPT_STATUS: unresolved
% ENTRY_KIND: research_attempt
\documentclass[11pt]{article}
\usepackage[margin=1in]{geometry}
\usepackage{amsmath,amssymb,amsthm,hyperref}
\newtheorem{proposition}{Proposition}
\newtheorem{lemma}{Lemma}
\title{Dynamic competition, innovation and mergers: a research attempt}
\author{GPT-6 (Codex)}
\date{9 October 2026}
\begin{document}
\maketitle
\noindent\textbf{Status: unresolved.} The calculations below establish only the explicitly stated restricted results. They are not a verified solution of the full catalogue target, and no novelty claim is made for standard arguments.

\section{The target and a deliberately explicit test model}
The original target asks for policy effects on innovation and welfare in a stochastic industry, allowing research, entry, acquisitions and equilibrium selection. A first useful test is whether a merger can have either innovation sign even in a uniquely solved subcase. If so, a universal sign claim requires restrictions absent from the broad formulation.

Consider two firms with research effort $a_i\in[0,A]$. Each period's expected quality improvement is $a_1+a_2$; effort is a calibrated intensity, not a patent count. Period profit is
\[
 \pi_i=b a_i+h a_j-g a_i a_j-\frac c2a_i^2+K_i,
 \quad b>0,\quad g\geq0,\quad c>2g.
\]
The private return $b$ may be the present value of a fixed lagged innovation payoff. The cross return $h$ captures a positive complement or negative business-stealing effect; $g$ represents diminishing joint returns. Constants $K_i$ do not affect research incentives. Repeat this one-state game with discount factor $\delta\in(0,1)$ and independent yearly innovations. Markov strategies are stationary, so current efforts solve the displayed stage optimization. This is a restrictive stochastic-game submodel, not a claim to model accumulated technology or entry.

\section{A sign threshold under unique behavior}
Before merger each first-order condition is $b-ga_j-ca_i=0$. Assuming the resulting effort lies strictly in $(0,A)$, the unique symmetric equilibrium is
\[
 a^N=\frac b{c+g}.
\]
Uniqueness follows because the joint best-response map has contraction modulus at most $g/c<1$, including clipping at the bounds. After merger a single owner maximizes $\pi_1+\pi_2$, giving
\[
 (b+h)-2ga_j-ca_i=0,\qquad a^M=\frac{b+h}{c+2g},
\]
provided $0<b+h<(c+2g)A$. The Hessian of joint profit has diagonal $-c$ and off-diagonal $-2g$, so it is negative definite when $c>2g$; the solution is uniquely optimal.

\begin{proposition}
In this interior class merger raises expected innovation if and only if
\[
 h(c+g)>bg.
\]
The discounted infinite-horizon innovation change is
$2(a^M-a^N)/(1-\delta)$.
\end{proposition}
\begin{proof}
Subtract the effort formulas:
\[
 a^M-a^N=\frac{h(c+g)-bg}{(c+2g)(c+g)}.
\]
The denominator is positive. Expected improvement is twice effort and the discounted sum is geometric.
\end{proof}
For $c=3$, $g=1$, $b=1$ and $A=1$, unmerged effort equals $1/4$. With $h=0$, merged effort equals $1/5$, so innovation falls. With $h=1$, it equals $2/5$, so innovation rises. Every displayed interior and concavity restriction holds. Thus the sign reversal does not require multiple equilibria, a change of innovation definition, or an ill-behaved optimization problem.

\section{What observed research incentives fail to identify}
The unmerged effort $b/(c+g)$ is independent of $h$. Firms' own optimality equations therefore do not identify the effect they impose on one another. If only efforts and net profits at the baseline equilibrium are observed, changing $h$ can be offset by changing $K_i$ to preserve those profit observations. The two examples can thus have identical baseline investment and profit data but opposite merger effects. Off-equilibrium variation in rival innovation, credible profit measurement and exclusion restrictions are needed to distinguish them.

A randomized change in rival research costs can, in principle, move $a_j$ without directly changing firm $i$'s return technology. Observing firm $i$'s profit response then identifies $h-ga_i$ after accounting for its own effort response. This statement requires the exclusion, the same innovation units and support for the induced variation; an industry-wide technology shock would not generally satisfy it. Merger policy instruments also need to avoid direct effects on financing and product demand if they are used to isolate this channel alone.

\section{Innovation and welfare have distinct signs}
Define real social benefit per expected unit of improvement as $v>0$, and count real research cost once. Suppose product-market surplus apart from innovation is $S_z$ under regime $z$, and a merger has real transaction cost $F\geq0$. The per-period innovation component of welfare at symmetric effort is
\[
 W(a)=2va-ca^2.
\]
Hence
\[
 W(a^M)-W(a^N)
 =(a^M-a^N)\{2v-c(a^M+a^N)\}.
\]
Even an innovation increase can reduce this welfare component if effort was already excessive relative to social return. Product-price effects $S_M-S_N$ and the transaction cost provide additional terms. Acquisition payments are transfers between owners and are not $F$. With unequal welfare weights or foreign ownership, those transfers may affect distributional welfare, but they still are not destroyed resources.

This calculation also shows why a patent count is insufficient: the conversion from counts to expected quality gain and its social value is a separate measurement equation. A merger can alter filing strategy without changing $a_1+a_2$.

\section{Selection-robust extensions and their limits}
For a genuine dynamic model at parameter $\theta$, define $V_z(\theta,e)$ as discounted innovation under equilibrium $e\in\mathcal E_z(\theta)$, using a common initial-state distribution. If equilibrium existence is established and
\[
 \inf_{\theta\in\Theta}\inf_{e\in\mathcal E_1(\theta)}V_1(\theta,e)
 >\sup_{\theta\in\Theta}\sup_{e\in\mathcal E_0(\theta)}V_0(\theta,e),
\]
then the innovation sign is positive for every permitted parameter and selection. This sufficient condition is stronger than necessary because it compares different parameter values across sides. The sharper test minimizes $V_1(\theta,e_1)-V_0(\theta,e_0)$ over a common $\theta$ and the actual restrictions linking selections. Computing only one converged equilibrium cannot certify either test.

For a fixed policy and a fixed equilibrium strategy, if a candidate value function has Bellman residual at most $\epsilon$ in sup norm, contraction bounds its value error by $\epsilon/(1-\delta)$. This verifies evaluation conditional on that strategy; it does not prove the strategy is an equilibrium or that other equilibria have been enumerated. Entry, ownership changes and privately known research productivity make those missing steps substantial.

\section{Completion Estimate}
31\%

This is a post hoc, heuristic estimate for the full catalogue target.
The attempt solves a unique research-effort merger benchmark with an exact innovation sign threshold and observationally identical opposite merger effects, and separates innovation from welfare. It does not identify dynamic industry outcomes with accumulating technologies, entry, acquisitions, private information, and sign robustness across equilibrium selections.

\section{Status and source check}
The effort comparison, counterexample pair, nonidentification observation and welfare accounting are complete in the announced submodel. The original industry-level dynamic identification and sign-robustness problem remains unresolved.

Richard Ericson and Ariel Pakes (1995), \emph{Markov-Perfect Industry Dynamics: A Framework for Empirical Work}, Review of Economic Studies 62(1), 53--82; publisher abstract inspected:
\url{https://academic.oup.com/restud/article-abstract/62/1/53/1568000}.
Ariel Pakes (2000), \emph{A Framework for Applied Dynamic Analysis in I.O.}, NBER WP 8024, primary abstract inspected for the computational framework and its limits:
\url{https://www.nber.org/papers/w8024}.

\end{document}
